We endow it with a @@M0@@-@@M1@@-algebra structure as follows: we observed above that the image of the functor @@M2@@ is in the subgroup of formal series in @@M3@@ whose coefficients in nonzero degree are in @@M4@@. Using the canonical morphism: @@M5@@, which gives: @@M6@@, we finally deduce a morphism:
@@M7@@
which, by 1.12, defines a morphism of functors on @@M8@@:
@@M9@@
The universal polynomials defining this morphism are the coefficients of the @@M10@@; the @@M11@@-th coefficient is the isobaric component of weight @@M12@@ of @@M13@@ (a proof analogous to that of 1.15), and will be denoted by @@M14@@.

Finally, we extend the morphism thus constructed to @@M15@@ naturally, using the @@M16@@-structure of @@M17@@. We therefore obtain a pre-@@M18@@-ring structure on @@M19@@, defined by the corresponding operations @@M20@@, which satisfy:
@@M21@@
and which will be denoted by @@M22@@ instead of @@M23@@, since no confusion is possible.

To verify that this is a @@M24@@-ring structure, we verify the commutativity of (3.7.2); it is known for the elements of @@M25@@; for the elements of @@M26@@, it follows from that of (3.7.2) for @@M27@@ by passage to the quotient. The result follows at once.

\providecommand{\Ghat}{\widehat G}
\providecommand{\Aqhat}{\widehat A_{\mathbb Q}}
\providecommand{\Ch}{\operatorname{Ch}}
\providecommand{\Log}{\operatorname{Log}}
\providecommand{\id}{\operatorname{id}}
\providecommand{\gr}{\operatorname{gr}}
\providecommand{\Hom}{\operatorname{Hom}}
\providecommand{\epsi}{\varepsilon}
\providecommand{\clgr}{\operatorname{cl}}
\providecommand{\Fil}{\operatorname{Fil}}
\providecommand{\Gr}{\operatorname{Gr}}
\providecommand{\Sym}{\operatorname{Sym}}

% Local macros for Exposé V, §§6--7
\providecommand{\Ghat}{\widehat G}
\providecommand{\Aqhat}{\widehat A_{\mathbb Q}}
\providecommand{\Ch}{\operatorname{Ch}}
\providecommand{\Log}{\operatorname{Log}}
\renewcommand{\id}{\operatorname{id}}
\providecommand{\gr}{\operatorname{gr}}
\providecommand{\Hom}{\operatorname{Hom}}
\providecommand{\epsi}{\varepsilon}
\providecommand{\clgr}{\operatorname{cl}}
\providecommand{\Fil}{\operatorname{Fil}}
\providecommand{\Gr}{\operatorname{Gr}}
\providecommand{\Sym}{\operatorname{Sym}}

